\documentclass[10pt,a4paper]{amsart}
\usepackage[margin=19mm]{geometry}
\usepackage{amsmath,amssymb,mathtools}
\usepackage[scr=rsfs,cal=euler]{mathalfa}
\usepackage{xfrac}
\usepackage[unicode,hidelinks,bookmarks=false]{hyperref}
\newcommand{\ie}{i.e.\ }
\newcommand{\cf}{cf.\ }
\newcommand{\resp}{respectively\ }
\newcommand{\Subsection}{\subsection}
\providecommand{\nobreakdash}{\nobreak\hbox{-}}
\setlength{\emergencystretch}{2em}
\multlinegap0pt
\usepackage{bm}
\usepackage{bbm}
\usepackage{dsfont}
\usepackage{xspace}
\usepackage{cancel}
\usepackage{mathtools}
\usepackage[shortlabels]{enumitem}
\usepackage{amsthm}
\makeatletter
\@ifundefined{theorem}{\newtheorem{theorem}{Theorem}}{}
\@ifundefined{lemma}{\newtheorem{lemma}[theorem]{Lemma}}{}
\@ifundefined{proposition}{\newtheorem{proposition}[theorem]{Proposition}}{}
\makeatother
\newcommand{\ZZ}{\mathbb{Z}}
\title[Coordinate projections of $c$-vectors of cluster algebras fr]{Coordinate projections of $c$-vectors of cluster algebras from the annulus}
\author{Sarah B. Brodsky}
\date{Source: arXiv:2606.27523v3 (CC BY 4.0)}
\begin{document}
\maketitle

\noindent\emph{Source and licence.} Coordinate projections of $c$-vectors of cluster algebras from the annulus. Version arXiv:2606.27523v3 (CC BY 4.0), distributed under the Creative Commons Attribution 4.0 International licence, as recorded in the arXiv licence metadata for this exact version and shown on the abstract page https://arxiv.org/abs/2606.27523v3. Full rights notice: licenses/arxiv-2606.27523-CC-BY-4.0.txt.\par\smallskip

\noindent\textit{Editorial scope.} Counted unit: the Proposition whose source label is \texttt{prop:surface} in the article (source0.tex lines 1471--1496, source label \texttt{prop:surface}), delivered together with the article's own complete original proof of it (source0.tex lines 1498--1573, one \texttt{proof} environment of 76 lines). No mathematics is written, repaired or re-derived: statement and proof are the article's own text line for line. Required context retained uncounted: thm:main (lines 152--167); lem:defect (lines 387--402); thm:defect\_detection (lines 714--727). Every definition, display and label of those lines is retained unchanged. Omitted from this package and not redistributed: 1--151, 168--386, 403--713, 728--1470, 1497--1497, 1574--1621. Nothing inside a retained excerpt is omitted or altered: every retained body is delivered exactly as the article's lines are written, so no omission receipt of any kind accompanies this package. Citations: the retained excerpts carry no citation command at all, so no bibliography entry of the article is transcribed and no reference is redistributed. Reference adaptation: none was needed and none was made; every reference inside the delivered unit resolves inside it, so no unresolved reference or citation is produced. Printed slips kept literal: no printed word, symbol, hypothesis or number of the counted statement or of its proof was altered. Layout and class: the article's own class is \texttt{amsart} with packages including amsmath, amssymb, amsthm, mathtools, enumitem, hyperref, tikz, booktabs; the wrapper is the lane's pinned \texttt{amsart} editorial wrapper at 10pt on A4, which restates the theorem-like environments the retained text needs and adds the article's own macro definitions that the retained text needs (\texttt{\textbackslash ZZ}), with extra wrapper packages (bm, bbm, dsfont, xspace, cancel, mathtools, enumitem). The delivered numbering is the unit's own. Assets: no retained span contains a figure environment, a graphics-inclusion command, a PDF-page inclusion, a table or a TikZ picture; no image file is copied into this package, the package declares no asset dependency and no image file is redistributed with it. Licence: version arXiv:2606.27523v3 of this article is distributed under the Creative Commons Attribution 4.0 International licence, as recorded in the arXiv licence metadata for this exact version and shown on the abstract page https://arxiv.org/abs/2606.27523v3. A full rights notice is delivered at licenses/arxiv-2606.27523-CC-BY-4.0.txt. Characters: every retained excerpt is pure ASCII, so the delivered engine, XeLaTeX with its default Latin Modern text font, reports no missing character.
\begin{theorem}[Projection-filling]
\label{thm:main}
Let $Q$ be a source-sink orientation of the $(n+1)$-cycle, with unique source
$s$ and unique sink $t$. Then:
\begin{enumerate}[label=\textup{(\arabic*)}]
\item every coordinate difference satisfies $|c_v - c_w| \leq 1$, so the
  band is $\{|c_v - c_w| \leq 1\}$;
\item for every pair $(v, w)$ other than $(s, t)$, the projection
  $\pi_{vw}$ surjects onto the full band;
\item for the source-sink pair, $\pi_{st}(C(Q))$ is the band with the
  diagonal $\{c_s = c_t\}$ removed except for a bounded segment; the
  diagonal carries only the regular part of $C(Q)$.
\end{enumerate}
In particular the source-sink diagonal $\{c_s = c_t\}$ is the unique
non-filled band line.
\end{theorem}

\begin{lemma}[Defect and transjective criterion]
\label{lem:defect}
For any acyclic orientation of the $(n+1)$-cycle, with sources $s_1, \dots,
s_k$ and sinks $t_1, \dots, t_k$, the defect form of $kQ$ is, up to sign,
\[
  \partial(x) \;=\; \sum_{i=1}^{k} x_{s_i} - \sum_{i=1}^{k} x_{t_i};
\]
this is a coordinate difference if and only if $k = 1$, that is, the source
and sink are unique, in which case $\partial(x) = x_s - x_t$. Under that
\textup{(}standing\textup{)} hypothesis a cyclic interval indicator
$\alpha_I$ is regular if and only
if $I$ contains both or neither of $\{s, t\}$, and transjective if and
only if $I$ contains exactly one of $\{s, t\}$. Moreover, if $\alpha_I$ is
transjective, then $\alpha_I + m\delta$ is a $c$-vector for every
$m \in \ZZ$, whereas the regular part of $C(Q)$ is finite.
\end{lemma}

\begin{theorem}[Projection-filling detects the defect]
\label{thm:defect_detection}
Let $Q$ be acyclic of affine type, and let $(v, w)$ be a banded pair with
$x_v - x_w = \pm\partial$ as linear forms. Then $\pi_{vw}$ fails to fill:
the diagonal $\{c_v - c_w = 0\}$ is the defect-zero locus, and it carries
only the finite regular part of $C(Q)$. For a source-sink orientation of type
$\widetilde{A}_n$, the defect is the coordinate difference
$\partial = (\cdot)_s - (\cdot)_t$, so the hypothesis holds for the
source-sink pair $(s, t)$; there $\partial$ takes only the values $-1, 0, +1$, the band is
$\{|c_s - c_t| \leq 1\}$, the three band lines $c_s - c_t = -1, 0, +1$
are exactly the preprojective, regular, and preinjective strata of
$C(Q)$, and the diagonal is the unique non-filled line. This is the
content of Theorem~\ref{thm:main}.
\end{theorem}

\begin{proposition}[Surface reading of the annulus theorem]
\label{prop:surface}
Let $Q$, $T$, and $c$ be as above, and let $t_c$ denote the Dehn twist
along $c$. Then:
\begin{enumerate}[label=\textup{(\arabic*)}]
\item an arc $\gamma$ is bridging if and only if $M_\gamma$ is
  transjective, and peripheral if and only if $M_\gamma$ is regular;
  equivalently
  \[
    i(\gamma, c) \;=\; \bigl|\partial(\underline{\dim}\, M_\gamma)\bigr|
    \;\in\; \{0, 1\};
  \]
\item $t_c$ fixes every peripheral arc; for a bridging arc $\gamma$, each
  coordinate of $k \mapsto \underline{\dim}\, M_{t_c^{\,k}\gamma}$ is a
  $V$-shaped sequence of slopes $\mp 1$ \textup{(}possibly with a flat
  bottom\textup{)}, on each of the two arms of the orbit consecutive
  twists eventually add exactly $\delta$, and the two arms carry constant
  defect $+1$ and $-1$, respectively;
\item consequently Theorem~\ref{thm:main} reads topologically: the two
  off-diagonal lines of $\pi_{st}$ carry the bridging arcs, separated by
  the sign of the defect \textup{(}the two arms of the twist
  orbits\textup{)}, and along each arm the Dehn twist realizes the
  $\delta$-shift filling the line; the diagonal carries only the finitely
  many arcs disjoint from the core.
\end{enumerate}
\end{proposition}

\begin{proof}
We first prove the crossing estimates for peripheral arcs. A peripheral
arc $\gamma$ has both endpoints on one boundary component, say the outer
one. Cutting $C_{p, q}$ along $\gamma$ leaves the inner boundary in a
single complementary region, an annular region into which $c$ can be
isotoped; so $\gamma$ is disjoint from a representative of $c$, is fixed
by $t_c$ up to isotopy, and has $i(\gamma, c) = 0$. Moreover, for any
bridging arc $\beta$, the minimal crossing number $i(\gamma, \beta)$ is
$0$ or $1$, and it depends only on the outer endpoint of $\beta$: if that
endpoint lies in the complementary region of $\gamma$ not containing the
inner boundary, then $\beta$ must cross $\gamma$ exactly once to reach
the inner boundary, and otherwise $\beta$ can reach the inner boundary
inside the other region and crosses $\gamma$ not at all. Observe now that
the two fan-junction arcs $s = (0, 0)$ and $t = (p, 0)$ have the
\emph{same} outer endpoint \textup{(}$0$ mod $p$\textup{)} and the same
inner endpoint \textup{(}$0$ mod $q$\textup{)}, differing only by
winding. Hence $i(\gamma, s) = i(\gamma, t)$; i.e.,
$\partial(\underline{\dim}\, M_\gamma) = 0$, and $M_\gamma$ is regular.

Next, every regular exceptional module is realized by a peripheral arc.
By Lemma~\ref{lem:defect} the regular real Schur roots of $Q$ are the
$\alpha_I$ with $I$ a cyclic interval containing both or neither of
$\{s, t\}$; equivalently, they are the $\alpha_I$ and $\delta - \alpha_I$
with $I$ an interval avoiding $\{s, t\}$. Such an interval consists of
consecutive arcs of one fan, say the outer one, whose outer endpoints
avoid the junction point $0$. The peripheral arc cutting off a disk
containing exactly the outer endpoints of the arcs of $I$ crosses exactly
the arcs of $I$, once each, realizing $\alpha_I$; and the peripheral arc
enclosing the inner boundary, chosen so that its inner-containing region
meets the outer boundary in exactly those endpoints, crosses exactly the
arcs \emph{not} in $I$, once each, realizing $\delta - \alpha_I$. Since
$\gamma \mapsto M_\gamma$ is a bijection onto the exceptional modules,
the peripheral arcs carry regular modules, and every regular exceptional
module is carried by a peripheral arc, the bridging arcs carry exactly
the transjective modules. Finally, a bridging arc meets the core, since
$c$ separates the two boundary components, and a representative crossing
once gives $i(\gamma, c) = 1$; combined with $|\partial| = 1$ on
transjective and $0$ on regular classes of $\widetilde{A}_n$
\textup{(}Lemma~\ref{lem:defect}\textup{)}, this proves the displayed
equality of (1).

For (2), the twist acts on a bridging arc by $(a, b) \mapsto (a + p, b)$,
which equals $(a, b - q)$ modulo the deck action. Writing $u = a - a_i$
and $w = b - b_i$ for the endpoint data relative to $\tau_i = (a_i,
b_i)$, the displayed crossing formula gives
\[
  i\bigl(t_c^{\,k}\gamma, \tau_i\bigr)
  \;=\; \#\Bigl\{\, j \in \ZZ :
  j \text{ lies strictly between } \tfrac{w}{q} \text{ and }
  k + \tfrac{u}{p} \,\Bigr\},
\]
which, as a function of $k$, is a $V$-shaped sequence with slopes
$\mp 1$, possibly with a flat bottom. For $k$ beyond the largest of the
thresholds $\frac{w}{q} - \frac{u}{p}$ over $i$, every coordinate
increases by exactly $1$ per twist; i.e., the dimension vector increases
by exactly $\delta$, and symmetrically in the other direction. For the
defect along the arms, observe that $s$ and $t$ have the same $w$ and
that their thresholds differ by exactly $1$. For $k$ large positive both
thresholds are exceeded, and the two counts differ by the number of
integers in the half-open unit interval
$[\,k + \frac{a - p}{p},\, k + \frac{a}{p}\,)$, which is one; for $k$
large negative the same window lies on the other side of $\frac{b}{q}$
and the difference is $-1$. Since $\partial$ is unchanged by adding
$\delta$, the defect is constant on each arm, equal to $+1$ on one and
$-1$ on the other, as claimed.

Part (3) is now a translation. The lines $c_s - c_t = \pm 1$ are the
transjective classes \textup{(}Theorem~\ref{thm:defect_detection}\textup{)},
which by (1) are the bridging arcs, split by the sign of $\partial$; by
(2) the two signs are carried by the two arms of the twist orbits, along
which consecutive twists add exactly $\delta$, so the Dehn twist realizes
the $\delta$-shift filling each off-diagonal line
\textup{(}Theorem~\ref{thm:main}\textup{)}. The diagonal is the regular
part, which by (1) consists of the peripheral arcs; these are disjoint
from the core and finite in number, as desired.
\end{proof}

\end{document}
